\section{Ablation Study: The Role-Complementarity
Principle}\label{sec:ablation}

The proposed AttentionHub-v2 is not introduced as a stand-alone design; it is
the conclusion of a systematic ablation. The benchmark of
Section~\ref{sec:results} shows \emph{that} the proposed model leads on the
subtype task, but not \emph{why} its hub is built from a Triplet~\arrowto{}SE
cascade rather than from any of the many attention configurations the
literature would licence. This section answers that question. It reports the
seven v1 ablation cells plus the no-attention control
(Section~\ref{subsec:ablcells}), isolates a quantitative regularity --- the
$98.36\,\%$ ceiling --- and abstracts it into a \emph{role-complementarity
principle} (Section~\ref{subsec:principle}). A confirmatory negative result
(Section~\ref{subsec:emaneg}) stress-tests the principle, after which the
principle is shown to force the v2 design (Section~\ref{subsec:v1tov2}) and the
resulting model is compared against the strongest baselines
(Section~\ref{subsec:bestv2}).

\subsection{Single-Branch and Pairwise Ablations}\label{subsec:ablcells}

The AttentionHub-v1 ablation is a code-level branch switch: each ablation key
maps to a subset of the three v1 branches $\{\text{BAM}, \text{Triplet},
\text{KAN}\}$, with the empty subset falling back to the donor's original
Block-4 (the canonical no-attention control). Table~\ref{tab:ablation} reports
the test-set accuracy of the eight resulting cells, sorted by ascending binary
accuracy. Every cell is trained under the identical baseline recipe used by the
nine baseline backbones, so the cells are directly comparable to one another
and to the benchmark of Section~\ref{sec:results}.

\begin{table}[!t]
\caption{AttentionHub-v1 ablation on the Custom EfficientNet~V2 backbone, under
the matched baseline recipe. Cells are sorted by ascending binary accuracy.
BAM is a mixed spatial-and-channel module, Triplet a cross-dimensional spatial
module, and KAN a purely channel-wise module. Role abbreviations: sp = spatial,
ch = channel, cd-sp = cross-dimensional spatial; $\cup$ denotes the union of
roles the active branches collectively exercise.}
\label{tab:ablation}
\centering
\setlength{\tabcolsep}{3.5pt}
\resizebox{\columnwidth}{!}{%
\begin{tabular}{l d{4} d{4} d{3} d{3} l}
\hline
\multicolumn{1}{l}{Variant} &
\multicolumn{1}{c}{Bin Acc} &
\multicolumn{1}{c}{Sub Acc} &
\multicolumn{1}{c}{Par. (M)} &
\multicolumn{1}{c}{GFLOPs} &
\multicolumn{1}{l}{Role(s)} \\
\hline
\texttt{bam\_triplet}  & 0.9825 & 0.9836 & 4.790 & 0.493 & sp+ch $\cup$ cd-sp \\
\texttt{kan}           & 0.9842 & 0.9897 & 4.777 & 0.490 & channel \\
\texttt{none}          & 0.9860 & 0.9885 & 5.700 & 0.636 & MBConv+SE \\
\texttt{triplet}       & 0.9895 & 0.9939 & 4.778 & 0.491 & cd-spatial \\
\texttt{bam}           & 0.9906 & 0.9909 & 4.779 & 0.491 & spatial+ch \\
\texttt{bam\_kan}      & 0.9906 & 0.9915 & 4.789 & 0.493 & sp+ch $\cup$ ch \\
\texttt{full}          & 0.9906 & 0.9921 & 4.800 & 0.495 & all three \\
\texttt{triplet\_kan}  & 0.9912 & 0.9945 & 4.788 & 0.493 & cd-sp + ch \\
\hline
\end{tabular}%
}
\end{table}

Three findings emerge from Table~\ref{tab:ablation}, illustrated by the bar
chart of Figure~\ref{fig:ablbars}.

\textbf{Finding~1: attention helps over the no-attention control on the
subtype task.} The donor Block-4 baseline (\texttt{none}) reaches $98.85\,\%$
subtype accuracy. Every Triplet-containing variant matches or exceeds this
figure with the single exception of \texttt{bam\_triplet}. The single-branch
\texttt{triplet} cell improves subtype accuracy to $99.39\,\%$ at \emph{fewer}
parameters than the no-attention control ($4.78$~M versus $5.70$~M), because
the donor's MBConv+SE block is heavier than a $1\times1$ channel-reduction
projection followed by a near-parameter-free Triplet gate. Attention here is
not merely additive --- it is additive \emph{and} cheaper.

\textbf{Finding~2: the \texttt{triplet\_kan} pair outperforms the full parallel
triple.} Removing BAM from the v1 hub \emph{improves} both heads:
binary accuracy rises from $99.06\,\%$ to $99.12\,\%$ and subtype accuracy from
$99.21\,\%$ to $99.45\,\%$. The full triple, although still strong, is not the
best v1 configuration --- a result obscured in the original proposal because
the full triple was never ablation-tested before publication. This finding is
the first concrete sign that adding a third branch can subtract performance.

\textbf{Finding~3: pairing two spatial-role modules collapses subtype accuracy
to a sub-baseline ceiling.} The \texttt{bam\_triplet} cell reaches only
$98.36\,\%$ subtype accuracy --- well below the no-attention control
($98.85\,\%$) and well below either constituent module evaluated singly
(\texttt{bam}: $99.09\,\%$; \texttt{triplet}: $99.39\,\%$). Both BAM and Triplet
exercise a spatial role (Section~\ref{subsec:principle}), and their combination
regresses \emph{beneath} the no-attention baseline rather than merely failing to
help. By contrast, \texttt{bam\_kan} --- which pairs BAM with the purely
channel-wise KAN --- reaches $99.15\,\%$ subtype, comfortably above the control.
The collapse is therefore specific to the spatial-role collision, not a generic
consequence of adding a second module. Section~\ref{subsec:emaneg} shows that
this same $98.36\,\%$ ceiling reappears in an independent configuration
(Triplet+EMA), confirming that the value is the signature of a spatial-role
conflict rather than a \texttt{bam\_triplet}-specific artefact.

\begin{figure}[!t]
\centering
\includegraphics[width=\columnwidth]{fig05_ablation_bars_with_ceiling}
\caption{Binary and subtype accuracy across the eight AttentionHub-v1 ablation
cells and the proposed v2 cascade. The dashed line marks the $98.36\,\%$
role-conflict ceiling: the spatial-spatial pair \texttt{bam\_triplet} lands
exactly on it, while every role-complementary configuration --- including
\texttt{bam\_kan} --- sits above it.}
\label{fig:ablbars}
\end{figure}

\subsection{The Role-Complementarity Principle}\label{subsec:principle}

Re-examining Table~\ref{tab:ablation} not by which modules a cell contains but
by \emph{what functional role each module exercises} resolves Finding~3 into a
clean pattern. Recall the role taxonomy of Section~\ref{sec:prelim}: BAM is a
mixed spatial-and-channel module, Triplet a cross-dimensional spatial module,
KAN a purely channel-wise module, and EMA a multi-scale spatial-and-channel
module. With modules so classified, the regularity in Table~\ref{tab:ablation}
can be stated as a principle.

\begin{quote}
\textbf{Role-complementarity principle.} \emph{Two attention modules combined
within a single hub should occupy disjoint functional roles. When both modules
exercise a spatial role, the hub regresses to a fixed sub-baseline accuracy
ceiling on this dataset; when one module is spatial and the other purely
channel-wise, the pair is complementary and improves over either module alone.}
\end{quote}

The principle is empirically supported, cell by cell, by
Table~\ref{tab:ablation}:

\begin{itemize}
\item \texttt{bam\_triplet} --- BAM (mixed spatial+channel) and Triplet
(cross-dimensional spatial) \textbf{both exercise a spatial role}. Subtype
$=98.36\,\%$ ($\times$).
\item \texttt{bam\_kan} --- BAM (mixed spatial+channel) and KAN (purely channel)
\textbf{do not collide on the spatial axis}; only one module is spatial.
Subtype $=99.15\,\%$ ($\checkmark$), above the no-attention control.
\item \texttt{triplet\_kan} --- Triplet (cross-dimensional spatial) and KAN
(purely channel) cover \textbf{disjoint roles}. Subtype $=99.45\,\%$
($\checkmark$).
\end{itemize}

Figure~\ref{fig:rolematrix} arranges the same evidence as a module-pair matrix:
the diagonal holds the single-module results and the off-diagonal holds the
pairs, with the spatial-spatial \texttt{bam\_triplet} cell appearing as the lone
sub-baseline entry. The $98.36\,\%$ ceiling recurs across \texttt{bam\_triplet}
and the independent Triplet+EMA configuration of
Section~\ref{subsec:emaneg} --- two spatial-role collisions reached from
different module sets --- which is the empirical signature of spatial-role
conflict rather than of a single bad cell. Because that number sits
\emph{below} the no-attention baseline of $98.85\,\%$, the failure mode is not
simply that the second module is uninformative: the role conflict
\emph{actively degrades} the representation that Stage~5 receives. A
redundant-but-harmless second module would, at worst, leave accuracy at the
single-module level; landing beneath the no-attention control is a stronger and
qualitatively different outcome.

\begin{figure}[!t]
\centering
\includegraphics[width=\columnwidth]{fig05b_role_complementarity_matrix}
\caption{Role-complementarity matrix. The diagonal reports single-module
subtype accuracy and the off-diagonal reports module-pair subtype accuracy.
The spatial-spatial pair BAM with Triplet is highlighted as the role-conflict
cell, regressing to the $98.36\,\%$ ceiling; the disjoint-role pair Triplet
with KAN reaches $99.45\,\%$, and the mixed pair BAM with KAN --- only one
spatial module --- reaches $99.15\,\%$, above the no-attention control.}
\label{fig:rolematrix}
\end{figure}

\subsection{Confirming the Principle: The v2-EMA Negative
Result}\label{subsec:emaneg}

A single failing cell within the v1 grid (\texttt{bam\_triplet}) invites two
competing explanations: that any second module degrades the hub, or that the
failure is specific to this one module pair. To separate them we performed a
non-ablation control: starting from \texttt{triplet\_kan}, we replaced KAN with
EMA (Efficient Multi-scale Attention, a multi-scale spatial-and-channel module)
and reorganised the hub into a sequential cascade --- the v2-EMA configuration.
Because EMA exercises a spatial role, the role-complementarity principle
predicts that v2-EMA will regress.

The prediction held. The v2-EMA configuration reached $98.60\,\%$ binary and
$98.36\,\%$ subtype accuracy --- inside the very same $98.36\,\%$ ceiling
produced by the \texttt{bam\_triplet} v1 cell. We report this as a documented
negative result. It independently corroborates the role-complementarity finding
from outside the v1 grid, and it rules out the competing explanation that
``any pair will do once a second module is added'': pairs do well when their
roles are disjoint and regress to the ceiling when their roles collide,
regardless of whether the colliding module is BAM or EMA. The $98.36\,\%$
ceiling is thus reached twice, from two disjoint module sets --- BAM+Triplet
within the v1 grid and Triplet+EMA outside it --- and in both cases by a pair of
spatial-role modules. This convergence is the empirical signature of
spatial-role conflict, and it rules out the competing explanation that any
second module degrades the hub: \texttt{bam\_kan} and \texttt{triplet\_kan},
whose second module is purely channel-wise, both improve on the no-attention
control.

\subsection{From v1 to v2: Sequential Triplet~\arrowto{}SE}\label{subsec:v1tov2}

The role-complementarity principle constrains the design of v2 directly. The
best v1 cell, \texttt{triplet\_kan}, pairs Triplet (cross-dimensional spatial)
with KAN (purely channel-wise), and the principle dictates that any improvement
over \texttt{triplet\_kan} must preserve that disjoint-role pairing: the
channel partner may be swapped for a different purely channel-wise module, and
the topology may be reorganised, but a spatial second module is forbidden. We
select Squeeze-and-Excitation (SE)~\cite{hu2018se} as the channel partner for
two reasons.

\begin{enumerate}
\item SE provides the same channel-recalibration function as KAN without the
small-data overfitting risk that KAN's learnable spline coefficients introduce.
On a dataset of this scale ($\approx5\,500$ training images across seven
classes) the per-channel B-spline parameters of KAN can begin to memorise
channel-specific idiosyncrasies; SE's two-layer $1\times1$ bottleneck, with far
fewer free parameters, is much harder to over-fit.
\item SE has been benchmarked extensively within CBAM and related
sequential-attention frameworks~\cite{woo2018cbam}, providing a
citation-grounded precedent for \emph{sequential} --- rather than parallel ---
attention composition.
\end{enumerate}

The topology is changed from parallel to sequential, following the CBAM
convention~\cite{woo2018cbam}: Triplet first refines spatial attention, and SE
then refines the channel response of the already spatially attended features.
No LayerScale and no per-module residual are used --- both attentions are
multiplicative gates that preserve input information without requiring an
explicit skip connection.

\textbf{Result.} Under the identical baseline recipe used by every v1 ablation
cell, the v2 cascade reaches $99.06\,\%$ binary and $99.51\,\%$ subtype
accuracy, improving both heads relative to the best v1 cell (\texttt{triplet\_kan}:
$99.12\,\%$~/~$99.45\,\%$) and the original v1 proposed model (\texttt{full}:
$99.06\,\%$~/~$99.21\,\%$). The $0.06$ percentage-point gain on the subtype
head over \texttt{triplet\_kan} represents a ten-image correction on a
$1\,646$-sample test set and is within single-run noise; we therefore present
v2 as the principled design, not as a statistically distinguishable
improvement over \texttt{triplet\_kan}. Figure~\ref{fig:progression} traces the
full design journey --- from the no-attention control, through the full v1
triple, to the best v1 cell, and finally to the v2 cascade --- with subtype
accuracy rising monotonically along the path while the binary head stays inside
its noise band.

\begin{figure}[!t]
\centering
\includegraphics[width=\columnwidth]{fig05d_v1_to_v2_progression}
\caption{AttentionHub design journey: no-attention control
($\rightarrow$) full v1 triple ($\rightarrow$) best v1 cell
\texttt{triplet\_kan} ($\rightarrow$) the proposed v2 Triplet~\arrowto{}SE
cascade. Subtype accuracy rises monotonically along the principled path; the
binary head remains within single-run noise throughout.}
\label{fig:progression}
\end{figure}

This entire chain --- ablation, principle, then v2 design --- is the central
methodological contribution of the paper. It replaces the ``stack attention
modules and report the best'' pattern common in the medical-imaging literature
with a measurement-driven design rule that explains, rather than merely
exhibits, the final architecture.

\subsection{Best v2 Configuration versus Strong Baselines}\label{subsec:bestv2}

Combining the conclusions of Sections~\ref{sec:results} and \ref{sec:ablation},
the proposed model under the role-complementarity principle is the smallest and
cheapest architecture in the entire study while delivering the highest subtype
accuracy among the nine baselines and a binary accuracy tied with the leader
within noise. Table~\ref{tab:proposed_vs_base} distils the head-to-head
comparison against the two strongest baselines.

\begin{table}[!t]
\caption{Digest of Tables~\ref{tab:benchmark} and \ref{tab:efficiency} for the
proposed model and the two strongest baselines. Bold marks the best
value in each row.}
\label{tab:proposed_vs_base}
\centering
\resizebox{\columnwidth}{!}{%
\begin{tabular}{l d{4} d{4} d{4}}
\hline
\multicolumn{1}{l}{Metric} &
\multicolumn{1}{c}{EffNet V2-B2} &
\multicolumn{1}{c}{Inception V3} &
\multicolumn{1}{c}{Custom V2 (v2)} \\
\hline
Binary accuracy  & \multicolumn{1}{c}{\textbf{0.9936}} & 0.9930 & 0.9906 \\
Subtype accuracy & 0.9933 & 0.9921 & \multicolumn{1}{c}{\textbf{0.9951}} \\
Parameters (M)   & 10.00 & 23.85 & \multicolumn{1}{c}{\textbf{4.79}} \\
GFLOPs           & 1.100 & 2.838 & \multicolumn{1}{c}{\textbf{0.493}} \\
Size (MB)        & 39.17 & 91.47 & \multicolumn{1}{c}{\textbf{18.94}} \\
GPU peak (MB)    & 72.47 & 132.18 & \multicolumn{1}{c}{\textbf{52.27}} \\
\hline
\end{tabular}%
}
\end{table}

Table~\ref{tab:proposed_vs_base} makes the trade-off explicit. The proposed
model concedes $0.30$ percentage points of binary accuracy to
EfficientNet~V2-B2 --- a margin already established as within single-run noise
in Section~\ref{subsec:quant} --- and in exchange leads on subtype accuracy and
on every efficiency metric by a wide margin: half the parameters of
EfficientNet~V2-B2 and a fifth of those of Inception~V3, under a quarter of
Inception~V3's GFLOPs, and the smallest on-disk and GPU footprint of the three.
The ablation-driven Triplet~\arrowto{}SE cascade is therefore not a marginal
re-tuning of an existing design; it is the principled endpoint of a
measurement-driven search, and it dominates the strongest baselines on the axes
that matter for a parameter-efficient classifier.
